\chapter{MARCO METODOLÓGICO}
\label{cap:metodo}

La metodología se organiza en tres unidades que se ejecutan en secuencia y que
cubren los objetivos específicos del proyecto: la preparación de la nube de
puntos, la extracción de descriptores y su fusión, y el entrenamiento,
calibración y evaluación del modelo. A ellas se añaden la delimitación del
dominio del objeto de estudio y el tratamiento de los riesgos. A continuación se
describe el diseño general del sistema y, después, cada una de las unidades.

% =====================================================================
\section{Diseño General del Sistema}
\label{sec:diseno}

El diseño del sistema y el significado de cada módulo, de cada fórmula y de
cada parámetro se describen en la Sección~\ref{sec:diseno-teorico} y se resumen
en la Figura~\ref{fig:diseno} y en la Tabla~\ref{tab:parametros}. Esta sección
no repite esa descripción: fija las propiedades de ejecución que el diagrama no
expresa, es decir, qué corre en paralelo, con qué frecuencia y con qué formato
de intercambio.

Las dos ramas están en paralelo porque ninguna consume la salida de la otra:
parten de la misma nube preparada y solo se encuentran en el Módulo~4. Esa
independencia permite ejecutar el Módulo~3 una sola vez, guardar su salida y
reutilizarla en todas las épocas del Módulo~2.

\subsection{Estructura y Contrato de los Módulos}
\label{sec:contrato}

Cada conexión entre módulos transporta un objeto con una forma determinada. La
Tabla~\ref{tab:contrato} declara ese contrato, de modo que cualquier
incompatibilidad de dimensión o de resolución aparezca antes de implementar.

\begin{table}[H]
\centering
\setstretch{1}\small
\setlength{\tabcolsep}{4pt}
\caption{Entrada, salida y régimen de cada módulo. Elaboración propia.}
\label{tab:contrato}
\begin{tabular}{@{}p{1.5cm}p{3.1cm}p{3.1cm}p{1.9cm}p{2.3cm}@{}}
\toprule
\textbf{Módulo} & \textbf{Entrada} & \textbf{Salida} & \textbf{Régimen} & \textbf{Frecuencia} \\
\midrule
1 & Nube TLS cruda, $N$ puntos & 5 secciones voxelizadas: coord., color, normal & Sin modelo & Una vez \\
2 & Sección voxelizada & $f_{\mathrm{sem}}$, 96 dim.\ por punto, rejilla dispersa & Entrenado & Cada época \\
3 & Sección voxelizada, 9 canales & $f_{\mathrm{ssl}}$, 1088 dim.\ por punto, rejilla serializada & Congelado & Una vez por sección \\
4 & $f_{\mathrm{sem}}$ y $f_{\mathrm{ssl}}$ alineados & Prototipos, hipergrafo y pseudo-etiquetas & Entrenado & Cada época \\
5 & Pseudo-etiquetas y verdad de terreno & Etiqueta por punto y mIoU por clase & Evaluación & Por corrida \\
\bottomrule
\end{tabular}
\end{table}

Las salidas de los Módulos~2 y~3 no viven en la misma estructura espacial. El
Módulo~2 trabaja sobre la rejilla dispersa de Minkowski y el Módulo~3 sobre su
propia rejilla serializada. Conciliarlas exige la propagación inversa del
submuestreo y un mapeo por índice inverso, ambos almacenados junto con las
características para que el Módulo~4 reciba los dos descriptores alineados punto
a punto.

% =====================================================================
\section{Unidad 1: Preparación de la Nube de Puntos}
\label{sec:unidad1}

La primera unidad corresponde al Módulo~1 y parte de la nube TLS cruda de la
Catedral de Lima entregada por la Facultad de Ingeniería Civil. Su propósito es
dejar los datos en una forma que las dos ramas de extracción puedan consumir sin
que ninguna etapa posterior requiera la nube completa en memoria.

La nube se recibe con al menos cuatro clases base anotadas en la sección de
referencia, requisito mínimo de la tarea NCD. Sobre ella se aplican tres
operaciones: el teselado en cinco secciones espaciales disjuntas (nave, coro,
ábside, torres y fachada), la voxelización a 5\,cm, que reduce el número de
puntos y homogeniza la densidad, y la normalización de coordenadas y atributos.
La salida de la unidad son las cinco secciones $S_1\ldots S_5$, cada una con
coordenadas, color y normal, que constituyen la unidad indivisible de trabajo
del proyecto: ninguna sección depende de otra, de modo que todas pueden
procesarse de forma independiente.

% =====================================================================
\section{Unidad 2: Extracción de Descriptores y Fusión}
\label{sec:unidad2}

La segunda unidad abarca los Módulos~2, 3 y~4. Las dos ramas de extracción
producen, sobre la misma sección preparada, los dos descriptores que el punto de
fusión concatena.

\subsection{Rama Semántica}

La rama semántica es la red troncal supervisada de HyperNCD, una
MinkowskiUNet-34C que opera sobre la rejilla dispersa y entrega
$f_{\mathrm{sem}}$, de 96 canales por punto. Es un módulo entrenado, de modo que
se ejecuta en cada época.

\subsection{Rama Auto-supervisada Congelada}
\label{sec:rama-ssl}

La rama auto-supervisada es el codificador PTv3 preentrenado de Sonata, usado
con pesos congelados y solo en inferencia. Entrega $f_{\mathrm{ssl}}$, de 1088
dimensiones por punto, y se ejecuta una sola vez por sección: su salida se
escribe en disco y se reutiliza en todas las épocas del entrenamiento.

La ejecución de esta rama produce un conjunto acotado de observables que el
diseño declara por adelantado y que la corrida debe confirmar
(Tabla~\ref{tab:salidas-p2}). Todos se registran en la misma corrida que genera
el almacenamiento intermedio, de modo que no exigen una ejecución adicional.

\begin{table}[H]
\centering
\setstretch{1}\small
\caption{Observables de la inferencia congelada del codificador. Elaboración
propia.}
\label{tab:salidas-p2}
\begin{tabular}{@{}p{3.9cm}p{4.3cm}p{4.3cm}@{}}
\toprule
\textbf{Observable} & \textbf{Para qué sirve} & \textbf{Valor esperado} \\
\midrule
Dimensión de $f_{\mathrm{ssl}}$ & Fija la entrada del Módulo~4 & 1088 \\
Puntos retenidos tras voxelizar & Determina el tamaño del almacenamiento y la resolución efectiva & Por medir \\
VRAM pico & Verifica la cota analítica de memoria & Cota estimada de 4{,}2\,GB \\
Tiempo de inferencia por sección & Sustituye a la estimación de costo & Por medir \\
Rendimiento en puntos por segundo & Permite extrapolar a otras secciones & Por medir \\
Tamaño del archivo de \emph{cache} & Dimensiona el almacenamiento del proyecto & Por medir \\
\bottomrule
\end{tabular}
\end{table}

Dos de estos observables condicionan decisiones de diseño. El tamaño de vóxel
efectivo del codificador no coincide necesariamente con el declarado para la
rejilla dispersa, porque la rutina de preparación de Sonata aplica su propio
submuestreo; verificar cuál es el valor aplicado determina el número de puntos
retenidos, la memoria y la resolución de salida. La disponibilidad de
FlashAttention determina si la matriz de atención se materializa de forma
explícita y, con ello, si el pico de memoria aparece o no.

\subsection{Punto de Fusión}
\label{sec:fusion}

El aporte consiste en conservar íntegro el razonamiento por hipergrafo y
sustituir la rama geométrica manual por la representación auto-supervisada,
según la Ec.~\eqref{eq:hibrido} del marco teórico. Aguas abajo de la fusión nada
cambia: la afinidad entre prototipos sigue combinando componente geométrica y
semántica según la Ec.~\eqref{eq:afinidad}, y cada prototipo conserva sus $M=8$
vecinos al construir el hipergrafo. Esa es la condición que hace interpretable
la ablación.

% =====================================================================
\section{Unidad 3: Entrenamiento, Calibración y Evaluación}
\label{sec:unidad3}

La tercera unidad describe cómo se ejecuta el sistema en el tiempo, con qué
procedimiento se obtienen las mediciones y bajo qué protocolo se repiten.

\subsection{Flujo de Ejecución}
\label{sec:diagrama-flujo}

\begin{figure}[H]
\centering
\begin{tikzpicture}[
  font=\scriptsize,
  paso/.style={draw, rounded corners, minimum width=5.6cm, minimum height=0.62cm,
               align=center, fill=blue!12},
  cong/.style={draw, rounded corners, minimum width=5.6cm, minimum height=0.62cm,
               align=center, fill=purple!14},
  disco/.style={draw, minimum width=5.6cm, minimum height=0.55cm,
                align=center, fill=black!8},
  dec/.style={draw, diamond, aspect=2.8, inner sep=1pt, align=center,
              fill=yellow!22},
  dat/.style={draw, rounded corners, minimum width=5.6cm, minimum height=0.5cm,
              align=center, fill=orange!18},
  fl/.style={-{Latex[length=1.8mm]}, thick},
  ret/.style={-{Latex[length=1.8mm]}, thick, black!60, rounded corners=2mm}
]
\node[dat] (p0) {\textbf{P0.} Nube TLS de la Catedral de Lima};
\node[paso, below=2.2mm of p0] (p1)
  {\textbf{P1.} Teselado, voxelización y normalización\\\emph{salida:} 5 secciones $S_1\ldots S_5$};
\node[cong, below=2.2mm of p1] (p2)
  {\textbf{P2.} Inferencia congelada de Sonata sobre $S_j$\\\emph{entorno CUDA 12.4}, una vez por sección};
\node[disco, below=2.2mm of p2] (cache)
  {\textbf{Frontera de disco:} \emph{cache} de $f_{\mathrm{ssl}}$ en \texttt{.npz}};
\node[paso, below=2.2mm of cache] (p3)
  {\textbf{P3.} Fusión y entrenamiento de HyperNCD sobre $S_j$\\\emph{entorno CUDA 11.4}, en cada época};
\node[dec, below=3.4mm of p3] (d1) {¿3 repeticiones\\completas?};
\node[dec, below=3.4mm of d1] (d2) {¿quedan\\secciones?};
\node[paso, below=3.4mm of d2] (p4)
  {\textbf{P4.} Agregación de métricas y análisis de dominancia};
\node[dat, below=2.2mm of p4] (p5)
  {\textbf{P5.} $\Delta$mIoU, brecha cerrada $G$ y ablación};

\foreach \a/\b in {p0/p1, p1/p2, p2/cache, cache/p3, p3/d1}
  \draw[fl] (\a) -- (\b);
\draw[fl] (d1) -- node[left=0.6mm, font=\tiny] {sí} (d2);
\draw[fl] (d2) -- node[left=0.6mm, font=\tiny] {no} (p4);
\draw[fl] (p4) -- (p5);

\coordinate (car) at ($(p2.east)+(0.8,0)$);
\coordinate (ciz) at ($(p1.west)+(-0.8,0)$);

\draw[ret] (d1.east) -- node[above, pos=0.7, font=\tiny, text=black!70] {no}
  (d1.east -| car) -- (car) -- (p2.east);
\node[font=\tiny, text=black!70, rotate=-90]
  at ($(car)!0.5!(d1.east -| car)+(0.3,0)$) {bucle interno};

\draw[ret] (d2.west) -- node[above, pos=0.7, font=\tiny, text=black!70] {sí}
  (d2.west -| ciz) -- (ciz) -- (p1.west);
\node[font=\tiny, text=black!70, rotate=90]
  at ($(ciz)!0.5!(d2.west -| ciz)+(-0.3,0)$) {bucle externo};

\node[below=3mm of p5, align=center, font=\tiny, text=black!70]
  {\tikz\node[draw,rounded corners,fill=purple!14,minimum height=2.4mm,minimum width=3.4mm]{};\ etapa congelada\quad
   \tikz\node[draw,rounded corners,fill=blue!12,minimum height=2.4mm,minimum width=3.4mm]{};\ etapa entrenable\quad
   \tikz\node[draw,fill=black!8,minimum height=2.4mm,minimum width=3.4mm]{};\ persistencia en disco\quad
   \tikz\node[draw,diamond,aspect=1.6,fill=yellow!22,minimum height=2.4mm,minimum width=3.4mm,inner sep=0.3pt]{};\ decisión};
\end{tikzpicture}
\caption{Flujo de ejecución del sistema. Elaboración propia.}
\label{fig:flujo}
\end{figure}

El flujo muestra tres elementos que la arquitectura no representa. El primero es
la frontera de disco: PTv3 requiere CUDA 12.4 y MinkowskiEngine la versión 11.x,
que no coexisten en un mismo proceso, por lo que $f_{\mathrm{ssl}}$ se escribe en
un archivo \texttt{.npz} y se lee desde el otro entorno. El segundo es la
asimetría de frecuencia: P2 se ejecuta una vez por sección y P3 en cada época,
lo que justifica el almacenamiento intermedio. El tercero son los dos bucles, el
de repeticiones y el de recorrido de secciones.

La diferencia con la Figura~\ref{fig:diseno} es de naturaleza. La
Figura~\ref{fig:flujo} es un diagrama de flujo porque describe una secuencia
temporal con decisiones y bucles; el diagrama de diseño describe componentes y
dependencias sin comprometerse con el orden de ejecución.
Si todo fuera secuencial bastaría el primero, pero los Módulos~2 y~3 son ramas
concurrentes, y esa concurrencia es una propiedad estructural que el flujo no
expresa.

\subsection{Procedimiento Experimental}
\label{sec:procedimiento}

El procedimiento se ejecuta en seis pasos reproducibles, correspondientes a los
rótulos P0 a P5 de la Figura~\ref{fig:flujo}. Los dos primeros son insumo de la
revisión crítica y fijan la línea base; los restantes corresponden a los
objetivos específicos:

\begin{enumerate}\setlength{\itemsep}{2pt}
\item \textbf{P0. Recepción.} Nube TLS con al menos cuatro clases base anotadas
en la sección de referencia, requisito mínimo de la tarea NCD.
\item \textbf{P1. Preparación y línea base.} Teselado en cinco secciones
disjuntas, voxelización a 5\,cm y normalización; ejecución de HyperNCD sin
modificaciones con tres semillas distintas para fijar el punto de partida contra
el que se compara el híbrido. El cambio de esquema de normalización fue ensayado
y descartado como resultado negativo.
\item \textbf{P2. Inferencia congelada (OE1).} Ejecución del codificador sobre
cada sección, con \emph{cache} de características y \emph{up-cast} para
conciliar la resolución del \emph{encoder} PTv3 con la rejilla de Minkowski, y
escritura de $f_{\mathrm{ssl}}$ en disco.
\item \textbf{P3. Fusión y entrenamiento (OE1).} Sustitución del descriptor
según la Ec.~\eqref{eq:hibrido} y verificación de funcionamiento de extremo a
extremo sobre una sección de control. Se comprueba que el mIoU de las clases
conocidas no caiga respecto de la línea base.
\item \textbf{P4. Calibración (OE2).} Búsqueda por grilla acotada sobre
$\alpha$, $M$ y $\epsilon$ en la sección de referencia; la configuración
ganadora se congela y se aplica sin reajuste al resto de secciones.
\item \textbf{P5. Análisis (OE3).} Cómputo de $\Delta$mIoU \emph{novel}
(Ec.~\eqref{eq:delta}), de la brecha cerrada $G$ (Ec.~\eqref{eq:G}) y del
análisis de dominancia (Ec.~\eqref{eq:dominante}), más la ablación de las tres
variantes del descriptor.
\end{enumerate}

\subsection{Protocolo de Repetición}
\label{sec:repeticiones}

Las secciones son independientes: ninguna necesita información de otra, de modo
que cada una constituye una medición separada y todas pueden ejecutarse de forma
simultánea.

Cada configuración se ejecuta tres veces sobre cada sección y se reporta la
media y la desviación entre repeticiones. La razón no es solo la variabilidad
del entrenamiento: el submuestreo que precede al codificador elige al azar el
punto representante de cada vóxel, por lo que dos ejecuciones sobre la misma
sección no producen la misma $f_{\mathrm{ssl}}$. Las tres repeticiones abarcan
por ello el procedimiento completo, P2 y P3, y no solo el entrenamiento.

\subsection{Trazabilidad de la Métrica}
\label{sec:trazabilidad}

Como el \emph{pipeline} es separable en dos etapas, la inferencia congelada de
Sonata y luego el razonamiento de HyperNCD, la métrica se reporta también en el
punto intermedio y no solo al final: \emph{linear probing} sobre el descriptor
antes de la formación de prototipos (Ec.~\eqref{eq:probe}) y mIoU después
(Ec.~\eqref{eq:miou}). Esa separación es la que permite atribuir una mejora a la
representación o al razonamiento, y no al sistema tratado como caja negra
(Sección~\ref{sec:metricas-hibrido}).

% =====================================================================
\section{Ventaja de Dominio del Objeto de Estudio}

Sonata se preentrenó sobre escenas interiores densas (ScanNet, S3DIS,
Structured3D, HM3D). En LiDAR urbano exterior su \emph{linear probing} queda por
debajo de PTv3 entrenado desde cero (62,0 frente a 69,1 mIoU en SemanticKITTI),
pero en interiores arquitectónicos densos alcanza 72,5 mIoU en ScanNet y 72,3 en
S3DIS Área~5, a la par de modelos supervisados completos. El levantamiento de una
catedral, estático, denso y de geometría arquitectónica, cae del lado
favorable de esa frontera, lo que reduce sustancialmente el principal riesgo de
transferencia de dominio.

% =====================================================================
\section{Riesgos y Mitigación}
\label{sec:riesgos}

\paragraph*{(i) Disponibilidad de la nube anotada.} NCD exige al menos las clases
conocidas etiquetadas. Se asume la anotación manual de una sección de referencia
con $\geq$4 clases base; si el levantamiento llega sin etiquetas, esa anotación
se vuelve tarea previa del cronograma y condiciona el inicio de OE1.

\paragraph*{(ii) Acceso y configuración del clúster.} La configuración de GPU de
Khipu ha cambiado respecto de la documentada previamente. Por ello el proyecto no
estima tiempos de calendario: verifica empíricamente el hardware asignado
(Tabla~\ref{tab:hardware}) y fija a partir de esa verificación el tamaño de
ventana y el esquema de \emph{cache}. El seccionamiento de la catedral es el
mecanismo que sostiene el respaldo: ninguna etapa requiere la nube completa en
memoria.

\paragraph*{(iii) Sustitución frente a concatenación (Plan~B).} Si el reemplazo
directo degrada el rendimiento, la variante es concatenar
$F_i=[f_{\mathrm{sem}};f_{\mathrm{geo}};f_{\mathrm{ssl}}]$, dejando que el módulo
de prototipos pondere las tres fuentes. Cualquier resultado, sea mejora, paridad o
degradación, documentado con ablación es reportable.

% =====================================================================
\section{Conclusión del Marco Metodológico}

Las tres unidades encadenan la preparación de la nube, la extracción de los dos
descriptores y su evaluación en un procedimiento que aísla una sola variable: el
descriptor que entra al módulo de prototipos. El contrato de la
Tabla~\ref{tab:contrato} fija las formas que cada módulo entrega, el flujo de la
Figura~\ref{fig:flujo} ordena las etapas en el tiempo y el protocolo de
repetición acota la variabilidad de las mediciones. De este modo, cualquier
diferencia observada en las métricas del Capítulo siguiente queda atribuida al
reemplazo de $f_{\mathrm{geo}}$ por $f_{\mathrm{ssl}}$ y no a la configuración
experimental.
